\section{Evaluasi}
\label{sec:evaluasi-bab3}

\subsection{Soal Pilihan Ganda}

\begin{enumerate}
    \item Bilangan biner adalah bilangan yang berbasis \dots
          \begin{enumerate}
              \item 2 dengan digit 0 dan 1
              \item 8 dengan digit 0 sampai 7
              \item 10 dengan digit 0 sampai 9
              \item 16 dengan digit 0 sampai F
          \end{enumerate}

    \item Nilai desimal dari \bin{1010} adalah \dots
          \begin{enumerate}
              \item 8 \item 10 \item 12 \item 16
          \end{enumerate}

    \item Nilai desimal dari \hex{F3} adalah \dots
          \begin{enumerate}
              \item 153 \item 243 \item 253 \item 263
          \end{enumerate}

    \item Pola 8 bit \bin{11111111} sebagai bilangan bertanda
          bernilai \dots
          \begin{enumerate}
              \item 255 \item 127 \item $-1$ \item $-128$
          \end{enumerate}

    \item Penulisan heksadesimal yang benar pada assembler adalah
          \dots
          \begin{enumerate}
              \item \code{A12H} \item \code{0A12H}
              \item \code{A12} \item \code{0A12D}
          \end{enumerate}
\end{enumerate}

\subsection{Soal Uraian}

\begin{enumerate}
    \item Ubah 200 desimal menjadi biner, oktal, dan heksadesimal.
          Tuliskan langkah pembagian dan pengelompokan bit yang Anda
          gunakan.

    \item Jelaskan perbedaan bilangan bertanda dan tidak bertanda pada
          8 bit. Tentukan bentuk \eng{two's complement} dari $-37$ dan
          periksa hasilnya dengan menghitung nilai tidak bertandanya.

    \item Pola \hex{F3} dapat berarti 243 atau $-13$. Jelaskan mengapa
          hal itu mungkin, dan sebutkan peran \eng{Sign Flag} serta
          jenis instruksi dalam menentukan tafsirnya.

    \item Tulislah potongan segmen data yang mendefinisikan 156 dalam
          bentuk desimal, biner, dan heksadesimal, ditambah karakter
          \code{'Z'}. Prediksikan isi memori yang tampil pada DEBUG,
          lalu jelaskan setiap byte.
\end{enumerate}

\subsection{Kisi-kisi Penilaian}

\begin{table}[H]
    \centering
    \caption{Kisi-kisi evaluasi Bab~\ref{chap:pengantar-assembly}.}
    \label{tab:kisi-bab3}
    \small
    \begin{tabular}{@{}p{1.5cm}p{5.0cm}p{2.6cm}c@{}}
        \toprule
        \textbf{Sub-CPMK} & \textbf{Indikator yang diukur} &
        \textbf{Bentuk soal} & \textbf{Bobot} \\
        \midrule
        1.1 & Menjelaskan tingkat bahasa dan peran assembly &
        Pilihan ganda 1 & 15\% \\
        \addlinespace
        1.1 & Mengonversi bilangan antar basis &
        Pilihan ganda 2, 3; Uraian 1 & 25\% \\
        \addlinespace
        1.1 & Membedakan bilangan bertanda dan tidak bertanda &
        Pilihan ganda 4; Uraian 2, 3 & 30\% \\
        \addlinespace
        1.1 & Menulis literal dan menafsirkan isi memori &
        Pilihan ganda 5; Uraian 4 & 30\% \\
        \bottomrule
    \end{tabular}
\end{table}

\begin{checklistbox}
\vspace{1.1em}
\textbf{Checklist Pencapaian Kompetensi
Bab~\ref{chap:pengantar-assembly}}

Seorang mahasiswa dinyatakan kompeten pada Sub-CPMK 1.1 apabila:
\begin{itemize}
    \item[\checked] dapat menjelaskan kedudukan bahasa assembly di
          antara bahasa mesin dan bahasa tingkat tinggi;
    \item[\checked] dapat mengonversi bilangan biner, oktal, desimal,
          dan heksadesimal dengan langkah yang benar;
    \item[\checked] dapat menuliskan literal bilangan yang diterima
          assembler, termasuk penggunaan nol di depan;
    \item[\checked] dapat membentuk dan menafsirkan
          \eng{two's complement};
    \item[\checked] dapat menjelaskan hubungan kode ASCII dengan data
          yang disimpan di memori;
    \item[\checked] dapat membaca dan menjelaskan hasil pemeriksaan
          memori dengan DEBUG.
\end{itemize}
\end{checklistbox}
